\section{Integration with the BX3 Framework}
\label{sec:rs-integration}

Recursive Spawning does not introduce its own accountability architecture — it
depends entirely on the three-layer structure defined in Section~\ref{sec:three-layer-model}.
Its correctness guarantees are derived from the structural properties that the
three layers enforce in concert.  This section specifies how Pillar 2 integrates
with each layer, completing the specification and preparing the ground for the
correctness properties in Section~\ref{sec:rs-correctness}.

% ---------------------------------------------------------------
\subsection{Purpose Layer as Accountability Anchor}
\label{sec:rs-purpose-anchor}

Every node $v$ in a BX3 spawn tree carries a human anchor $h(v)$ that is
established at node provisioning and is structurally non-modifiable for the
node's operational lifetime (Invariant~\ref{inv:human-anchor}).  The
\textit{Purpose} layer of $v$ is the architectural locus of this anchor: it
carries the principal's intent, defines the operating range $R(v)$, and is the
exclusive terminus for any escalated exception that the \textit{Bounds} layer
cannot resolve within that range.

The integration between Recursive Spawning and the \textit{Purpose} layer
operates at three distinct points:

\begin{enumerate}
\item[(a)] \textbf{Anchor inheritance at spawn.}  When parent node $p$ spawns child
$c$, the spawn directive encodes $h(c) = h(p)$.  The \textit{Purpose} layer of
$c$ inherits the same human anchor as its parent.  This is not a policy decision
made by the agent at $p$ — it is a structural constraint enforced by the
\fact{} layer's pre-commit hook before the child node is provisioned.  The
implication is that \textit{every node in the spawn tree, regardless of depth,
points to the same human principal} as its ultimate accountability terminus.

\item[(b)] \textbf{Escalation terminus.}  The Bailout Protocol traverses the
parent-pointer chain $\alpha(v), \alpha(\alpha(v)), \ldots$ until it reaches
$h(v)$.  Because $h(v)$ is always a human principal and never a machine actor,
the escalation path never terminates at an intermediate machine node.  The
\textit{Purpose} layer at each node is the last machine layer before the human
anchor — it holds the principal's intent, not a delegated subset of it.  This
structural property makes the upstream accountability guarantee non-vacuous: the
guarantee holds precisely because the terminus is defined at the \textit{Purpose}
layer, which cannot be replicated or substituted by any machine actor.

\item[(c)] \textbf{Scope adjudication.}  When the \textit{Bounds} layer detects
a delegation scope inheritance violation (Invariant~\ref{inv:scope-inheritance}),
the exception is escalated to the \textit{Purpose} layer of the parent node for
adjudication.  The \textit{Purpose} layer alone has the authority to redefine
$R(p)$ or to authorize an exception to the inheritance constraint on behalf of
the human principal.  No machine actor may grant this authorization.
\end{enumerate}

The consequence of (a)–(c) is that the \textit{Purpose} layer serves as a
double anchor: it anchors the spawn tree's accountability chain at the top
(human principal $\rightarrow$ root node's \textit{Purpose} layer), and it
anchors every escalation path at its terminus ($h(v)$, reached through the
parent-pointer chain).  The spawn tree is not a tree of autonomous agents — it
is a tree of accountability structures, each anchored to the same human principal.

% ---------------------------------------------------------------
\subsection{Bounds Engine: Depth and Scope Enforcement}
\label{sec:rs-bounds-depth}

The \textit{Bounds} layer enforces two deterministic pre-commit checks on every
spawn event: the Hierarchy Finiteness check and the Delegation Scope Inheritance
check.  Both are enforced as integer and set comparisons respectively — no
probabilistic reasoning, no model judgment.

\medskip
\noindent\textbf{Depth check.}  The depth-evaluation function
$\text{depth-ok}(c, p)$ computes $d(c) = d(p) + 1$ and compares the result
against the deployment constant $D_{\max}$.  When a node in state
\texttt{CAN\_SPAWN} receives an $e_{\text{spawn}}$ event, it transitions to
\texttt{SPAWN\_PENDING} and issues a depth-confirmation request to the
\textit{Bounds} layer.  If $d(c) \leq D_{\max}$, the \textit{Bounds} layer emits
$e_{\text{depth-ok}}$; if not, it emits $e_{\text{depth-violation}}$,
transitions the spawning state machine to \texttt{SPAWN\_LOCKED}, triggers the
Bailout Protocol, and records the violation in the Forensic Ledger.

\medskip
\noindent\textbf{Scope inheritance check.}  The \bounds{} layer evaluates
$R(c) \subset R(p)$ before every spawn event as a deterministic set comparison.
If the subset relation fails, the \bounds{} layer emits
$e_{\text{scope-violation}}$, transitions the spawning state machine to
\texttt{SPAWN\_LOCKED}, and triggers the Bailout Protocol.  The scope check
precedes the depth check in evaluation order, ensuring that the first detected
violation is the one recorded in the Forensic Ledger — which matters for
accountability attribution when the human anchor reviews the failure.

The critical property of this integration is that $D_{\max}$ is a deployment
parameter established at system initialization and is not negotiable at runtime.
No machine actor, including the \bounds{} layer itself, may raise $D_{\max}$ to
accommodate a non-compliant spawn.  The deployment designer sets $D_{\max}$ once,
based on the operational requirements of the deployment context — specifically,
the maximum acceptable escalation path length for the Bailout Protocol's Liveness
guarantee.  The \bounds{} layer's enforcement of the depth constraint ensures that
this contract between deployment designer and runtime system is never violated.

% ---------------------------------------------------------------
\subsection{Fact Layer: Forensic Ledger}
\label{sec:rs-fact-fledger}

The \fact{} layer provides two structural guarantees for recursive spawning:
immutable parent-pointer enforcement and complete, cryptographically chained event
recording.

\medskip
\noindent\textbf{Immutable parent pointer.}  The parent-pointer field
$\alpha(v)$ is a read-only property of the node's \fact{} layer worksheet.  Any
attempt to write to $\alpha(v)$ at runtime is intercepted by the pre-commit hook
and rejected.  The rejection event is recorded as a
\texttt{parent-pointer-write-denied}$(v, \alpha_{\text{old}}, \alpha_{\text{new}})$
entry, and the Bailout Protocol is triggered if the attempted write was
initiated by a machine actor.  This enforcement ensures that the accountability
chain is never redirected after node creation.

\medskip
\noindent\textbf{Cryptographically chained forensic record.}  Every spawn event,
every state transition in the spawning state machine, every depth- and
scope-confirmation result, and every human directive issued in response to a
spawn-related exception is recorded as a Forensic Ledger entry.  Each entry $f$
carries the SHA-256 hash of the preceding entry, forming an append-only chain
from the genesis entry to the most recent.  The cryptographic chaining ensures
that no entry can be modified, deleted, or inserted without breaking the chain —
a condition detectable by any auditor without access to the ledger's contents.

The Forensic Ledger entries relevant to recursive spawning are:

\begin{itemize}\itemsep=0pt
\item \texttt{spawn-proposed}$(p, c, R(c), d(c))$ — parent $p$ proposes child
$c$ with operating range $R(c)$ and depth $d(c)$;
\item \texttt{depth-confirmed}$(c, d(c))$ — \bounds{} layer confirmed $d(c) \leq D_{\max}$;
\item \texttt{depth-rejected}$(c, d(c), D_{\max})$ — depth check failed;
\item \texttt{scope-confirmed}$(c, R(c) \subset R(p))$ — \bounds{} layer confirmed
delegation scope inheritance;
\item \texttt{scope-rejected}$(c, R(c), R(p))$ — scope-inheritance check failed;
\item \texttt{child-provisioned}$(p, c, \alpha(c) = p, h(c) = h(p))$ — child
successfully created with immutable parent pointer and inherited human anchor;
\item \texttt{parent-pointer-write-denied}$(v, \alpha_{\text{old}}, \alpha_{\text{new}})$
— write to $\alpha(v)$ intercepted and rejected;
\item \texttt{spawn-locked}$(v, \text{violation-type})$ — spawning state machine
entered \texttt{SPAWN\_LOCKED} due to a constraint violation.
\end{itemize}

The \fact{} layer's pre-commit hook queries the Forensic Ledger to determine the
current state of the spawning state machine before committing any spawn event,
ensuring that the ledger is the authoritative source of truth for the runtime
state of the spawn tree.  The record is not a conventional audit log — it is an
enforcement mechanism that records state transitions before subsequent actions
are permitted.

% ---------------------------------------------------------------
\subsection{Unified Integration Property}
\label{sec:rs-unified-integration}

The three-layer integration yields a single structural property that is the
foundation of all downstream correctness guarantees:

\begin{invariant}[Spawn Tree Accountability Invariant]
\label{inv:spawn-accountability}
For every node $v$ in a BX3 spawn tree, the accountability chain
$\alpha^{(k)}(v) = h(v)$ with $k \leq D_{\max}$ is: (i) immutable — $\alpha(v)$
is never modified after node creation; (ii) complete — every spawn event and state
transition is recorded in the Forensic Ledger; (iii) terminating — the chain
terminates at a human principal $h(v)$, never at a machine actor; and (iv)
enforceable — any violation of (i), (ii), or (iii) triggers the Bailout Protocol
before the system state advances.
\end{invariant}

This invariant is enforced by the pre-commit hooks of all three layers operating
in concert.  The \purpose{} layer provides the human anchor; the \bounds{} layer
provides depth and scope enforcement; the \fact{} layer provides the immutable
parent pointer and the cryptographically chained forensic record.  None of the
three layers alone is sufficient.  It is their conjunction that enables the
correctness properties established in the next section.
